\documentclass{article}
\usepackage[T2A]{fontenc}
\usepackage[utf8]{inputenc}   
\usepackage[english, russian]{babel}

% Set page size and margins
% Replace `letterpaper' with`a4paper' for UK/EU standard size
\usepackage[a4paper,top=2cm,bottom=2cm,left=2cm,right=2cm,marginparwidth=1.75cm]{geometry}

\usepackage{amsmath}
\usepackage{graphicx}
\usepackage[colorlinks=true, allcolors=blue]{hyperref}
\usepackage{amsfonts}
\usepackage{amssymb}
% \usepackage[left=1cm,right=1cm,top=1cm,bottom=1cm]{geometry}
\usepackage{hyperref}
\usepackage{seqsplit}
\usepackage[dvipsnames]{xcolor}
\usepackage{enumitem}
\usepackage{algorithm}
\usepackage{algpseudocode}
\usepackage{algorithmicx}
\usepackage{mathalfa}
\usepackage{mathrsfs}
\usepackage{dsfont}
\usepackage{caption,subcaption}
\usepackage{wrapfig}
\usepackage[stable]{footmisc}
\usepackage{indentfirst}
\usepackage{rotating}
\usepackage{pdflscape}
\usepackage{listings}

\usepackage{minted}

\usepackage{MnSymbol,wasysym}

\title{Домашнее задание 2}
\author{Лиза Фоменко}
\date{}

\begin{document}

\maketitle



\section*{\textbf{Задание 1}} \textbf{\textit{Найдите $\Theta$-асимптотику функции $f(n) = \sum^n_{k=0} {n \choose k}$}}

Докажем по индукции с использованием следующих фактов:

\begin{enumerate}
  \item $\binom{n}{0} = \binom{n+1}{0} = 1$
  \item $\binom{n}{n} = \binom{n+1}{n+1} = 1$
  \item $\binom{n}{k} + \binom{n}{k-1} = \binom{n+1}{k}$; доказательство этого факта: $\binom{n}{k} + \binom{n}{k-1} = \frac{n!}{k!(n-k)!} + \frac{n!}{(k-1)!(n-k+1)!} = \frac{n!}{(k-1)!(n-k)!} * (\frac{1}{k} +\frac{1}{n-k+1}) =  \frac{n!}{(k-1)!(n-k)!} * \frac{n-k+1+k}{k*(n-k+1)} = \frac{n!}{(k-1)!(n-k)!} * \frac{n+1}{k*(n-k+1)} = \frac{(n+1)!}{k!(n-k+1)!} = \binom{n+1}{k}$
\end{enumerate}

\medskip
Далее воспользуемся индукцией по n.

\smallskip

\textbf{База индукции:} $2^1 = \binom{1}{0} + \binom{1}{1} = 1 + 1 = 2$

\smallskip


\textbf{Предположение индукции:} $2^n = \sum^n_{k=0} {n \choose k}$

\smallskip


\textbf{Шаг индукции:} $2^{n+1} = 2*2^n = 2*\sum^n_{k=0} {n \choose k} = 2*(\binom{n}{0} + \binom{n}{1} + \binom{n}{2} + ... + \binom{n}{n-2} + \binom{n}{n-1} + \binom{n}{n}) = 2*\binom{n}{0} + 2*\binom{n}{1} + 2*\binom{n}{2} + ... + 2*\binom{n}{n-2} + 2*\binom{n}{n-1} + 2*\binom{n}{n} = \binom{n}{0} + (\binom{n}{0} + \binom{n}{1}) + (\binom{n}{1} + \binom{n}{2}) + ... + (\binom{n}{n-2} + \binom{n}{n-1}) + (\binom{n}{n-1} + \binom{n}{n}) + \binom{n}{n} = \binom{n+1}{0} + \binom{n+1}{1} + \binom{n+1}{2} + ... + \binom{n+1}{n-2} + \binom{n+1}{n-1} + \binom{n+1}{n+1} = \sum^{n+1}_{k=0} {n+1 \choose k}$

\medskip

По мат индукции доказали, что $2^n = \sum^n_{k=0} {n \choose k}$. Отюсда воспользуемся так называемым "свойством рефлексивности" из Кормена: $f(n) = \Theta(f(n)) \Rightarrow \sum^n_{k=0} {n \choose k} = \Theta(2^n)$

\bigskip

\section*{\textbf{Задание 2}} \textbf{\textit{Дан набор из $n$ монет, одна из которых фальшивая и легче других, и чашечные весы, на которые можно положить две группы монет одинакового размера и узнать, какая из них тяжелее, или что они равны по весу. Придумайте алгоритм, докажите его корректность и оцените асимптотику.}}

оформлю словесный алгоритм в виде псевдокода :)

\begin{verbatim}
while True:
    if число монет в куче n mod 3 == 0:
        делим кучу на 3 равные части: куча1 = n/3,  куча2 = n/3,  куча3 = n/3 монет
    elif число монет в куче mod 3 == 1:
        делим кучу на 3 части: куча1 = n//3, куча2 = n//3 и куча3 = n//3 + 1 монет
    else:
        делим кучу на 3 части:  куча1 = n//3 + 1, куча2 = n//3 + 1, куча3 = n//3 монет
        
    взвешиваем кучу1 и кучу2 (всегда равны по кол-ву монет):
        if одинаковые по весу:
            куча = куча3 
            if число монет в куче == 1:
                return куча (нашли монетку)
        else:
            куча = min(куча1, куча2)
            if число монет в куче == 1:
                return куча (нашли монетку)
\end{verbatim}

\medskip

\textbf{Корректность:} на i-м шаге мы ищем $min(kucha_{i1}, kucha_{i2}, kucha_{i3})$ за 1 сравнение, где $kucha_{ik}$ - одна из трех частей $kucha_i$, которую мы приняли на вход нашего i-го шага. 

\textit{\textbf{Инвариант:}} на i шаге $kucha_{i} = min(kucha_{(i-1)1}, kucha_{(i-1)2}, kucha_{(i-1)3}$. 

\textit{\textbf{База индукции:}} на первом шаге мы делим исходную кучу на 3 части -- $kucha_{11}, kucha_{12}$ и $kucha_{13}$. За одно взвешивание $kucha_1$ и $kucha_2$ с равным количеством монеток (важно для работы алгоритма!) получим, что либо они одинаковы $\Rightarrow$ легкой монетки в них нет $\Rightarrow$ она в $куче_3$, либо одна из них будет легче $\Rightarrow$ легкая монетка в более легкой из них. Тогда $kucha_2 = min(kucha_{11}, kucha_{12}, kucha_{13}$. 

\textit{\textbf{Шаг индукции:}} $kucha_{i+1} = min(kucha_{i1}, kucha_{i2}, kucha_{i3} = min(kucha_i // 3, kucha_i // 3, kucha_i // 3) = min(min(kucha_{(i-1)1}, kucha_{(i-1)2}, kucha_{(i-1)3} // 3, min(kucha_{(i-1)1}, kucha_{(i-1)2}, kucha_{(i-1)3} // 3, min(kucha_{(i-1)1}, kucha_{(i-1)2}, kucha_{(i-1)3} // 3)$. (здесь показано целочисленное деление для простоты, но зависит от случая :) )

\medskip

\textbf{Ассимптотика:} На каждом шаге алгоритма мы:
\begin{enumerate}
    \item делим кучу на 3 части $\Theta(1)$
    \item взвешиваем 2 из них $\Theta(1)$
    \item выбираем самую легкую кучу $\Theta(1)$
\end{enumerate}

То есть каждый шаг алгоритма выполняется за $\Theta(1)$ (это, конечно, некоторого рода абстракция; скорее всего, в реализации кодом "взвешивание" выглядело бы как к-л сумма / лин операция, и ее нельзя было бы считать за $\Theta(1)$). Остается только лишь подсчитать, сколько шагов мы сделаем: с каждым шагом размер кучи уменьшается ~втрое, поэтому шагов будет порядка $\left\lceil \log_3n \right\rceil$ (пока не дойдем до кучек размера 1), где $n$ - общее количество монет. Тогда: $\left\lceil \log_3n \right\rceil * \Theta(1) = \Theta(\log_3 n)$.

\textbf{ответ:} ассимптотика алгоритма $\Theta(\log_3 n)$. Это оптимальный алгоритм.

\bigskip


\section*{\textbf{Задание 3}} \textbf{\textit{Найдите нижнюю оценку для задачи поиска фальшивой монеты, не предполагая ничего про алгоритм. Нужно найти нижнюю оценку сложности именно для задачи, а не для какого-либо конкретного алгоритма, решающего ее.}}

Пусть у нас есть дерево, где узлы - сравнения (взвешивания), а листья - исходы всех взвешиваний (условно, i лист = i-я монета фальшивая). С каждого узла мы имеем 3 возможных исхода: =, <, >. Тогда наше дерево - ~полное M-арное, где M=3 (тернарное, кажется), а количество листьев = n. Тогда высота такого дерева составляет $\left\lceil \log_3n \right\rceil$. Высота дерева в данном контексте -- это количество взвешиваний, которые могут довести нас от изначальной кучи до фальшивой монеты (само значение высоты при известном количестве листьев использую без доказательства, но оно аналогично бинарному дереву, пример в Кормене). 

Тогда для любой входной кучи из n монет мы можем за $\left\lceil \log_3n \right\rceil$ взвешиваний дойти до любого(!!) его листа. 

Если же высота  дерева будет меньше, чем  $\left\lceil \log_3n \right\rceil$, например $\left\lceil \log_3n \right\rceil$ - 1? Дерево остается тернарным, ~полным, и для него верно, что при высоте $\left\lceil \log_3n \right\rceil$ - 1 максимальное количество его листьев (= возможных исходов) будет равно $3^{\left\lceil \log_3n \right\rceil - 1} = n / 3$, то есть за такое количество взвешиваний у нас останется доступ только к n/3 возможным исходам: перемешав кучу, мы можем не попасть в нужный лист (нарушение детерменированности алгоритма).  Поэтому нжняя оценка алгоритма составляет $\Omega(\left\lceil \log_3n \right\rceil)$. 

Мой алгоритм работает за $\Theta(\log_3 n)$, те ялвяется оптимальным.
 

\section*{\textbf{Задание 4}} \textbf{\textit{Найдите $7^{13} \mod 167$ с помощью быстрого возведения в степень. Нужно привести последовательность умножений и промежуточные результаты. mod 167 - это остаток от деления на 167.}}

Будем использовать быстрое вовзведение в степень за $\log_2 n$ и следующий факт: (a*b) mod p = (a mod p) * (b mod p) mod p.

\begin{enumerate}
    \item 7 mod 167 = 7
    \item $7^2 $ mod 167 = (7 mod 167)$^2$ mod 167 = 7*7 mod 167 = 49
    \item $7^3$ mod 167 = ($7^2$ mod 167) * (7 nod 167) mod 167 = (49 * 7) mod 167 = 343 mod 167 = 9
    \item $7^6$ mod 167 = ($7^3$ mod 167)$^2$ mod 167 = 9*9 mod 167 = 81
    \item $7^{12}$ mod 167 = ($7^6$ mod 167)$^2$ mod 167 = 81*81 mod 167 = 6561 mod 167 = 48
    \item $7^{13}$ mod 167 = ($7^{12}$ mod 167) * (7 mod 167) mod 167 = (48 * 7) mod 167 = 2
\end{enumerate}

\textbf{ответ:} 2 


\section*{\textbf{Задание 5}}  \textbf{\textit{Рассмотрим последовательность $L_i$, которая похожа на числа Фибоначчи, но в обратную сторону. Она будет начинаться с двух положительных целых чисел $L_0 = a$ и $L_1 = b \leq a$, а следующие элементы будут определяться как $L_{i+1} = \max(L_i, L_{i-1}) - \min(L_i, L_{i-1})$. Определите, сколько ненулевых элементов может быть в такой последовательности, если она начинается с двух чисел, запись большего из которых состоит из $n$ бит. Как изменится ответ, если вместо вычитания брать остаток от деления нацело?}}


\textbf{1)} $L_{i+1} = \max(L_i, L_{i-1}) - \min(L_i, L_{i-1})$

Мне не очень понятен вопрос "сколько" элементов. Имеется в виду сколько \textit{различных} элементнов может быть в последовательности или сколько \textit{всего} ее членов являются ненулевыми?

Рассмотрим оба варианта постановки вопроса. Начнем с вопроса \textit{сколько всего} ненулевых членов в последовательности $L_i$, для этого рассмотрим следующий пример:

Пусть у нас конечное число ненулевых элементов в последовательности, тогда (так как она в целом убывающая) в какой-то момент, начиная с $L_k, \exists k \in \mathbb{N}$, что для $\forall n \geq k$ верно, что $L_n = 0$. Тогда рассмотрим $L_{k+1}$ и $L_k$. Мы значем, что $L_{k+1} = |L_k - L_{k-1}| = |0 - L_{k-1}| = |L_{k-1}| = 0 \Rightarrow L_{k-1} = 0$. Однако мы предположили, что последовтельность 0 начинается только с $L_k$ ее члена. Следуя такой логике, мы получим, что вся такая последоватльность состоит только из 0.

Рассмотрим случай, когда $L_k$ оказалось равна 0. Это значит, что $L_k = |L_{k-1} - L_{k-2}| = 0 \Rightarrow L_{k-1} = L_{k-2}$. Тогда два следующих члена: $L_{k+1} = |L_{k} - L_{k-1}| = |0 - L_{k-1}| = |L_{k-1}| = L_{k-2}; L_{k+2} = |L_{k+1} - L_{k}| = |L_{k-1} - 0| = |L_{k+1}| = |L_{k-1}| = L_{k-2}$. То есть когда мы получаем 0 в последовательности, она начинает выглядеть как $x, x, 0, x, x, 0, ...$ и так бесконечно.

То есть $L_i$ имеет бесконечное общее число ненулевых элементов. Разберемся с числом различных элементов в этой последовательности.

Мы знаем, что $a$ состоит из $n$ бит, то есть максимальное значение $a = 11...1$ ($n$ единиц). Также $a \geq b > 0$, то есть наша последовательность точно будет убывающей. Сколько всего существует $n$-битных чисел (то есть у нас $n$ позиций, на каждой из которых может стоять либо 1, либо 0)? Ответ: $2^n$ чисел. Тогда попробуем подобрать пример, в котором встретятся все $2^n$ чисел.

\begin{enumerate}
    \item $L_0 = a = 11...11$, n единиц
    \item $L_1 = b = 1 = 1$
    \item $L_2 = |L_0 - L_1| = |a-b| = |11...11 - 1| = 11...10$
    \item $L_3 = |L_1 - L_2| = |1 - 11...10| = 11...01$
    \item $L_4 = |L_2 - L_3| = |11...01 - 11...10| = 1$
    \item $L_5 = |L_3 - L_4| = |11...01 - 1| = 11..100$
    \item $L_6 = |L_4 - L_5| = |1 - 11...00| = 11..011$
    \item $L_7 = |L_5 - L_6| = |11..011 - 11..100| = 1$
    \item etc, $L_{3k+1} = 1, L_{3k-1} - L_{3k} = 1$
\end{enumerate}

\medskip

Как видно из приведенных выше первых членов и их общих формул, составленная таким образом $L_i$ будет содержать все $2^n$ чисел (ее конечная повторяющаяся бесконечная часть будет выглядеть как $1,1,0,1,1,0,...$. 

\textbf{Итого ответ:} разных ненулевых чисел - $2^n - 1$; всего ненулевых чисел - бесконечное количество (так как с к-л момента последовательность становится периодичной).


\bigskip

\textbf{2)} $L_{i+1} = \max(L_i, L_{i-1}) \mod \min(L_i, L_{i-1})$

Наша последовательность прервется, когда $L_k = 0$ для какого-либо k (так как мы работаем с положительными числами, то 0 будет минимальным в паре с любым другим числом, и мы в последовательности будем пытаться искать остаток от деления $L_{k-1}$ на $L_k = 0$, получим неопределенность). Тогда, дойдя до первого 0 в последовательности, она завершится. При этом наша последовательность всегда будет конечна, так как частное от деления всегда меньше, чем делитель, а мы начинаем с целого числа $a$.  

Здесь нам необходимо максимизировать число ненулевых элементов, то есть построить наиболее медленно убывающую последовательность: для этого необходимо получать наибольший возможный остаток, что (в свою очередь) будет при частном = 1.

Тогда рассмотрим $L_{i-1}, Li, L_{i+1}$: зная, что частное = 1, запишем выражение для $L_{i+1}$: $ L_{i+1} = L_i * 1 + L_{i-1} = L_i + L_{i-1} \Rightarrow $ мы получили последовательность Фиббоначчи, с тем лишь небольшим отличием, что начало истинной послед. Фиббоначчи выглядит как 0, 1, 1, 2, 3; а конец нашей последовательности - 3, 1, 0 (так любое число при делении на 1 дает остаток 0).

Наибольшая последовательность нужного вида $F_k, F_{k-1}, ..., 3, 2, 1, 0$, где $F_k$ - самое большое $n-$битное число Фиббоначчи.


Мне, к сожалению, не удалось вывести общий вид числа Фиббоначчи в двоичном коде, чтобы по данному $n$ общей формулой вывести $F_k$ (наибольшоее $n$-битное число Фиббоначчи). Замечу, что если по $n$ научиться вычислять $k$, то длина такой последовательности будет $k-2$ (так как получим последовательность ненулевых чисел от $F_k$ до $F_2 = 1$, $F_1=1$ не используется (см выше), а $F_0 = 0$).

\end{document}